% Part: first-order-logic
% Chapter: axiomatic-deduction
% Section: provability-quantifiers

% verification of properties of provability needed for maximally
% consistent sets in the completeness chapter.

\documentclass[../../../include/open-logic-section]{subfiles}

\begin{document}

\olfileid{fol}{axd}{qpr}

\olsection{\usetoken{S}{derivability} आणि संख्यापक}

\begin{explain}
  संपूर्णता प्रमेयासाठी या विभागात स्थापित केलेली~$\Proves$ संबंधी तथ्ये
  स्वयंसिद्धकीय निगमनातून निष्पन्न होतात, हेही आवश्यक आहे.
\end{explain}

\begin{thm}
\ollabel{thm:strong-generalization} जर $c$ हे $\Gamma$ किंवा $!A(x)$ मध्ये
न येणारे !!a{constant} असेल आणि $\Gamma \Proves !A(c)$, तर $\Gamma
\Proves \lforall[x][!A(x)]$.
\end{thm}

\begin{proof}
निगमन प्रमेयाने $\Gamma \Proves \ltrue \lif !A(c)$. $c$ हे
$\Gamma$ किंवा~$\top$ मध्ये येत नसल्यामुळे $\Gamma \Proves
\ltrue \lif \lforall[x][!A(x)]$ मिळते. त्यानंतर सत्य स्थिरांकाचे
स्वयंसिद्धक आणि विध्यात्मक अनुमान वापरून $\Gamma \Proves
\lforall[x][!A(x)]$ मिळते.
%READERNOTE{OLAXD-006: गोठवलेल्या इंग्रजी मजकुरात शेवटचे पाऊल पुन्हा निगमन प्रमेयाने मिळते असे म्हटले आहे. प्रत्यक्षात आधीच्या सशर्त निष्कर्षावर सत्य स्थिरांकाचे स्वयंसिद्धक आणि विध्यात्मक अनुमान लावल्यावर अपेक्षित निष्कर्ष मिळतो; मराठीत हे कारण स्पष्ट केले असून गोठवलेले इंग्रजी बाइट्स बदललेले नाहीत.}
\end{proof}

\begin{prop}
\ollabel{prop:provability-quantifiers}
\begin{tagenumerate}{prvEx,prvAll}
\tagitem{prvEx}{$!A(t) \Proves \lexists[x][!A(x)]$.}{}

\tagitem{prvAll}{$\lforall[x][!A(x)] \Proves !A(t)$.}{}
\end{tagenumerate}
\end{prop}

\begin{proof}
\begin{tagenumerate}{prvEx,prvAll}
\tagitem{prvEx}{\olref[qua]{ax:q2} आणि निगमन प्रमेय वापरून.}{}
\tagitem{prvAll}{\olref[qua]{ax:q1} आणि निगमन प्रमेय वापरून.}{}
\end{tagenumerate}
\end{proof}

\end{document}
